% Diplomatic French for printed p.218.
% Printed A in (5.6.11 ter), the extra citation period, and spectral are preserved.
\clearpage
\SourcePage{218}

\LayerTarget{tohoku:transition:degeneration-case}
\LayerTarget{tohoku:corollary:theorem-5.6.3-locally-free}
\textsc{Corollaire.}\quad \textit{Si le $G$-${\mathbf O}$-module $A$ est
localement isomorphe, en tant que ${\mathbf O}$-module, à ${\mathbf O}^n$, alors on a des
isomorphismes canoniques}

\LayerTarget{tohoku:equation:5.6.12}
\begin{equation}
\operatorname{Ext}_{{\mathbf O},G}^n(X;A,B)
  =H^n\bigl(X;G,\mathbf{Hom}_{\mathbf O}(A,B)\bigr).
\tag{5.6.12}
\end{equation}

\LayerTarget{tohoku:proof:theorem-5.6.3-locally-free-corollary}
Cela résulte de la suite spectrale $III$, puisque alors
$\mathbf{Ext}_{\mathbf O}^q(A,B)=0$ pour $q>0$ (prop. 4.2.3). En particulier, prenant
$A={\mathbf O}$, on trouve :

\LayerTarget{tohoku:corollary:theorem-5.6.3-structure-sheaf}
\textsc{Corollaire 2.}\quad \textit{On a, pour tout $G$-${\mathbf O}$-Module
$B$, des isomorphismes fonctoriels :}

\LayerTarget{tohoku:equation:5.6.14}
\begin{equation}
\operatorname{Ext}_{{\mathbf O},G}^n(X;{\mathbf O},B)=H^n(X;G,B).
\tag{5.6.14}
\end{equation}

\LayerTarget{tohoku:remarks:two-additional-spectral-sequences}
\textsc{Remarques.}\quad Signalons encore, dans deux cas
particuliers, deux autres suites spectrales aboutissant à
$\operatorname{Ext}_{{\mathbf O},G}^n(X;A,B)$.
\LayerTarget{tohoku:construction:constant-equivariant-crossed-product}
Supposons que ${\mathbf O}$ soit le $G$-faisceau d'anneaux constant défini par un anneau
$O$ où $G$ opère, et que $A$ soit le $G$-${\mathbf O}$-faisceau constant $\mathbf M$
défini par un $O$-module $M$ où $G$ opère de façon compatible avec ses
opérations sur $O$ : $g\cdot(\lambda m)=g(\lambda)g(m)$. Introduisons
l'anneau $U$ engendré par $O$ et $\mathbf Z(G)$ soumis aux relations de
commutation $g\lambda g^{-1}=g(\lambda)$, $M$ est donc un $U$-module.
Considérons le foncteur $h_{U,M}$ sur la catégorie des $U$-modules, défini par
$h_{U,M}(N)=\operatorname{Hom}_U(M,N)$, on a alors pour tout $G$-${\mathbf O}$-Module
$B$ :
$\operatorname{Hom}_{{\mathbf O},G}({\mathbf M},B)=\operatorname{Hom}_U(M,\Gamma_X(B))$, i.~e.
on a l'identité fonctorielle

\LayerTarget{tohoku:equation:5.6.5-bis}
\begin{equation}
h_{{\mathbf O},G,M}=h_{U,M}\Gamma_X.
\tag{5.6.5 bis}
\end{equation}
\LayerTarget{tohoku:spectral-sequence:iv}
On vérifie facilement de plus que $\Gamma_X$ transforme objets injectifs de
${\mathbf C}^{{\mathbf O}(G)}$ en des $U$-modules injectifs, de sorte que le théorème 2.4.1 donne
un foncteur spectral cohomologique $IV$ sur la catégorie ${\mathbf C}^{{\mathbf O}(G)}$,
aboutissant au foncteur gradué
$\bigl(\operatorname{Ext}_{{\mathbf O},G}^n(X;{\mathbf M},B)\bigr)$, et dont le terme initial est

\LayerTarget{tohoku:equation:5.6.11-bis}
\begin{equation}
IV_2^{p,q}(B)=\operatorname{Ext}_U^p\bigl(M,H^q(X,B)\bigr).
\tag{5.6.11 bis}
\end{equation}
\LayerTarget{tohoku:specialization:trivial-action-on-m}
Si on suppose de plus que $G$ opère trivialement sur $M$, alors on a aussi
$\operatorname{Hom}_{{\mathbf O},G}({\mathbf M},B)=\operatorname{Hom}_U(M,\Gamma_X(B))=\operatorname{Hom}_{{\mathbf O}}(M,\Gamma_X^G(B))$,
d'où la formule fonctorielle

\LayerTarget{tohoku:equation:5.6.5-ter}
\begin{equation}
h_{{\mathbf O},G,{\mathbf M}}=h_{O,M}\Gamma_X^G
\tag{5.6.5 ter}
\end{equation}
\LayerTarget{tohoku:spectral-sequence:v}
On peut encore montrer que $\Gamma_X^G$ transforme objets injectifs de
${\mathbf C}^{{\mathbf O}(G)}$ en $O$-modules injectifs (en montrant que si $N$ est un $U$-module
injectif, alors $N^G$ est un $O$-module injectif), et le théorème 2.4.1. donne
un cinquième foncteur spectral cohomologique, de même aboutissement que les
précédents, et dont le terme initial est

\LayerTarget{tohoku:equation:5.6.11-ter}
\begin{equation}
V_2^{p,q}(B)=\operatorname{Ext}_{\mathrm O}^p\bigl(M,H^q(X;G,A)\bigr)
\tag{5.6.11 ter}
\end{equation}
\LayerTarget{tohoku:comparison:spectral-sequences-iv-v}
Si $G$ est réduit à l'unité, les suites spectrales $IV$ et $V$ coïncident avec la
suite spectrale du théorème 4.3.1. Si $X$ est réduit à un point, la suite
spectrale $IV$ est triviale, mais la suite spectrale $V$ n'est pas triviale en
général, et ne se réduit pas aux précédentes\textsuperscript{12)}\SourceFootnote{12}{\LayerTarget{tohoku:footnote:12}Le terme initial de cette suite spectral est
(avec les notations de 11))
$\operatorname{Ext}_{\mathrm O}^p\bigl(A,H^q(G,B)\bigr)$.}.

\TohokuContinuation{219}
